\clearpage
\VerdierPhysicalPageAnchor{0131}%
\begingroup
\small
\setlength{\emergencystretch}{1.5em}%
\setlength{\abovedisplayskip}{2pt plus 1pt minus 1pt}%
\setlength{\belowdisplayskip}{2pt plus 1pt minus 1pt}%

\VerdierSemanticAnchor{ch2-2-1-5-continuation}
\noindent には、包含関手と両立する三角構造は高々一つしか存在しない。

\VerdierSemanticAnchor{ch2-2-1-5-criterion}
$\mathcal D'$ をシフトで安定な $\mathcal D$ の充満部分圏とする。
$\mathcal D'$ が三角構造をもち、$\mathcal D'$ が $\mathcal D$ の
充満三角部分圏となるための必要十分条件は、$\mathcal D'$ の任意の二対象
$X,Y$ と任意の射
$f\colon X\to Y$ に対し、$Z$ が $\mathcal D'$ の対象であるような
$\mathcal D$ の完全三角 $(X,Y,Z,f,g,h)$ が存在することである。

\VerdierSemanticAnchor{ch2-2-1-6}
\medskip
\noindent\textbf{2.1.6.}\quad
\VerdierSemanticAnchor{ch2-2-1-6-saturated}
$\mathcal D$ を三角圏とする。$\mathcal D$ の充満三角部分圏 $\mathcal D'$
が次の性質をもつとき、\emph{飽和}であるという。

\VerdierSemanticAnchor{ch2-2-1-6-direct-summand}
\par\smallskip
\noindent\makebox[0pt][l]{(2.1.6.1)}%
\begin{minipage}[t]{\linewidth}
  \centering\itshape
  $\mathcal D'$ の対象の $\mathcal D$ における任意の直和因子は、\\[-1pt]
  $\mathcal D'$ の対象と同型である。
\end{minipage}
\par\smallskip

\VerdierSemanticAnchor{ch2-2-1-6-direct-summand-definition}
\noindent （対象 $Y$ の直和因子とは、ある対象 $Z$ と同型
$X\oplus Z\xrightarrow{\sim}Y$ が存在するような対象 $X$ をいう。）

\VerdierSemanticAnchor{ch2-2-1-7}
\medskip
\noindent\textbf{2.1.7.}\quad
\VerdierSemanticAnchor{ch2-2-1-7-denominator-set}
$\mathcal D'$ を三角圏 $\mathcal D$ の充満三角部分圏とする。次で定義される
$\mathcal D$ の射の集合を $S(\mathcal D')$ と書く：

\VerdierSemanticAnchor{ch2-2-1-7-characterization}
\par\smallskip
\noindent\makebox[0pt][l]{(2.1.7.1)}%
\begin{minipage}[t]{\linewidth}
  \centering\itshape
  $f\colon X\to Y$ が $S(\mathcal D')$ に属するための必要十分条件は、\\[-1pt]
  $Z$ が $\mathcal D'$ の対象であるような $\mathcal D$ の完全三角
  $(X,Y,Z,f,g,h)$ が存在することである。
\end{minipage}
\par\smallskip

\VerdierSemanticAnchor{ch2-2-1-8}
\medskip
\noindent\textbf{命題 2.1.8.}\quad
\VerdierSemanticAnchor{ch2-2-1-8-proposition}
\emph{集合 $S(\mathcal D')$ (2.1.7.1) は $\mathcal D$ の乗法的系
(2.1.4) である。それが飽和であるための必要十分条件は、部分圏
$\mathcal D'$ が飽和であることである。}

\VerdierSemanticAnchor{ch2-2-1-8-proof-plan}
第一の主張を証明するには、(2.1.3) により、集合 $S(\mathcal D')$ が性質
SM1)、SM3)、SM5)、SM6) をもつことを示せばよい。以下
$S(\mathcal D')$ がこれらの性質をもつことを順に示し、最後に末尾の主張を
証明する。

\VerdierSemanticAnchor{ch2-2-1-9}
\medskip
\noindent\textbf{2.1.9.}\quad
\VerdierSemanticAnchor{ch2-2-1-9-sm1}
\emph{$S(\mathcal D')$ は性質 SM1) をもつ。}

\VerdierSemanticAnchor{ch2-2-1-9-identity}
まず、$\mathcal D$ の任意の対象 $X$ に対し、射 $\operatorname{id}_X$ が
$S(\mathcal D')$ に属することは明らかである。これは TRI) と、
$\mathcal D'$ の零対象が $\mathcal D$ の零対象でもあることから直ちに従う。

\VerdierSemanticAnchor{ch2-2-1-9-composition}
$S(\mathcal D')$ が合成について閉じていることは、公理 TRIV) と (1.2.4)
から直ちに従う。検証は読者に委ねる。

\VerdierSemanticAnchor{ch2-2-1-10}
\medskip
\noindent\textbf{2.1.10.}\quad
\VerdierSemanticAnchor{ch2-2-1-10-sm3}
\emph{$S(\mathcal D')$ は性質 SM3) をもつ。}

\VerdierSemanticAnchor{ch2-2-1-10-reduction}
まず、$s\in S(\mathcal D')$、$fs=0$ なる図式
$X\xrightarrow{s}Y\xrightarrow{f}Z$ が与えられたとき、
$t\in S(\mathcal D')$、$tf=0$ なる図式
$Y\xrightarrow{f}Z\xrightarrow{t}W$ が存在することを示せばよい。
そうすれば反対圏へ移ることで逆向きも得られ、その結果、

\VerdierSemanticAnchor{ch2-2-1-10-p132-boundary}
\endgroup
